\CUSTOMstruct{ЗАКЛЮЧЕНИЕ}

В ходе производственной, технологической практике выполнена проверка компетенций обучающегося путём работы над настоящими задачами в проекте: 
отработаны навыки постановки и согласования задач, работы в команде, написания и тестирования кода, внедрения нового функционала в уже имеющиеся системы, профилирования при нагрузке, сформирован отчёт о выполненных этапах.

В рамках разработки утилиты для сбора логов по событиям Camunda р
реализован REST-сервис, формирующий выгрузку истории процессов с фильтрацией по временному диапазону. 
Проведено профилирование нагрузки, подтвердившее неинтрузивность решения. 
Усвоены подходы к проектированию вспомогательных сервисов, не влияющих на производительность основной системы.

В рамках доработки получения движений токена по workflow система 
переведена с периодического опроса на событийную модель. 
Реализован плагин Camunda Engine, активируемый при старте, с поддержкой корректного восстановления после перезапуска. 
Получен практический опыт построения event-driven архитектур.

В рамках настройки движка Camunda под нагрузку оптимизированы Job
Executors для асинхронного выполнения external tasks, отключена 
запись истории в слой персистенции.

В рамках написания кеша для метаданных реализована библиотека, наполняемая при старте из базы данных и инвалидируемая по событиям. 
Подтверждено снижение числа обращений к базе данных. Усвоены принципы проектирования переиспользуемых библиотек с минимальными зависимостями и паттерн кеширования с событийной инвалидацией.

В результате практики освоены навыки работы с Camunda BPM: интеграция с интерфейсом исторических активностей,
разработка плагинов для системы бизнес-процессов,      
настройка отдельного модуля под
нагрузку. Приобретён опыт        
проектирования event-driven
систем на базе брокера сообщений,
реализации кеша с инвалидацией по событиям.